\chapter{Інформаційна технологія AURA/Ecliptix та експериментальна перевірка}
\label{ch:infotech}

% ============================================================
%  РОЗДІЛ 8 -- інтеграція методів у програмну інформаційну технологію.
%  Доводить інженерну реалізованість та інтеграцію, НЕ нову теорему.
%  Canonical backend: aura-platform.
% ============================================================

\section{Архітектура реалізації}
\label{sec:impl-arch}
% -- Rust AURA Core; iOS client; backend services; gateway.

\section{Інтерфейси міжмовної взаємодії (FFI)}
\label{sec:impl-ffi}
% -- OPAQUE FFI; protected protocol FFI; межі та контракти.

\section{Інтеграція каналу нагляду та межа приватності пошуку}
\label{sec:impl-integration}
% -- supervision relay integration; search privacy boundary.

\section{Керування наборами даних}
\label{sec:dataset-governance}
% -- curated/synthetic корпуси, побудовані за моделлю месенджерних діалогів;
%    governance, gold/mixed зрізи.

\section{Release-gated оцінювання та privacy-safe аудит}
\label{sec:release-eval}
% -- support-aware release gates; статуси PASS/FAIL/INSUFFICIENT_SUPPORT/BLOCKED;
%    privacy-safe audit records.

\section{Експериментальна перевірка}
\label{sec:experiments}
% -- Тести, бенчмарки, відтворюваність; порівняння з baselines де доречно.

\section{Індекс артефактів та відтворюваність}
\label{sec:artifact-index}
% -- Таблиця: repo / commit / метод-розділ / команда перевірки / очікуваний результат / межі.
% [Деталі -- Додаток А.]

\section{Обмеження та межі застосування}
\label{sec:limitations}
% -- Чесний перелік обмежень; не clinical validation; напрями подальших досліджень.
Окремо узагальнено обмеження технології, зокрема відсутність клінічної валідації safety-рішень та класичний характер схеми підпису, що визначає межі застосування й напрями подальших досліджень. \emph{[Розкрити обмеження.]}

\rozdilconcl{8}
\emph{[Висновки до розділу 8.]}
